\Def  Группой Галуа многочлена $f \in F[x]$ называется группа автоморфизмов, сохраняющих $F$, его поля разложения. Обозначение — $Gal f$

\Def Группа $G$ разрешимая, если $\exists n \in \N $: n-ый коммутант $G^{(n)}$ тривиален.

\Note Из теории групп: разрешимость $\Leftrightarrow $ в G $\exists$ цепочка $G = G_n \unrhd G_{n-1} \dots \unrhd G_{0} = e$: $\forall i$ фактор-группа $G_{i+1}\backslash G_i$ абелева.

\Th Для $f(x) \in F[x]$ $f(x) = 0$ разрешимо в радикалах $\Leftrightarrow$ Gal $f$ разрешима.

\Proof 
Пусть $L = L_0 \supset \dotsb \supset L_n = F$ "--- башня расширений из определения разрешимости в радикалах. Проверим, что $\Gal{f} = \Aut_{L_n}{L} \ge \dotsb \ge \Aut_{L_0}L = \{\id\}$ при добавлении нескольких промежуточных подгрупп образует цепочку нормальных подгрупп с абелевыми факторами.
		
		Для произвольного $i \in \{0, \dotsc, n - 1\}$ рассмотрим многочлен $x^{m_{i+1}} - a_{i+1} \in L_{i + 1}[x]$, полем разложения которого является $L_i$, а также $K_i$ "--- поле разложения многочлена $x^{m_{i+1}} - 1 \in L_{i + 1}[x]$, тогда $L_i = L_{i + 1}(\sqrt[\leftroot{2}\uproot{4}m_{i+1}]{a_{i+1}}, \xi_{m_{i+1}}) \supset K_i = L_{i + 1}(\xi_{m_{i+1}}) \supset L_{i + 1}$.
		
		Каждый автоморфизм $\phi \in \Aut_{L_{i+1}}K_i$ однозначно задается образом элемента $\xi_{m_{i+1}}$, а $\xi_{m_{i+1}}$ может переходить только в один из элементов вида $\xi_{m_{i+1}}^m$, где $m \in \N$ "--- такое, что $(m, {m_{i+1}}) = 1$, причем композиция автоморфизмов соответствует умножению степеней элемента $\xi_{m_{i+1}}$. Значит, $\Aut_{L_{i+1}}K_i \subset \Z_{m_{i+1}}^*$, поэтому группа $\Aut_{L_{i+1}}K_i$ "--- абелева.
		
		Аналогично, каждый автоморфизм $\phi \in \Aut_{K_i}L_i$ однозначно задается образом элемента $\sqrt[\leftroot{2}\uproot{4}m_{i+1}]{a_{i+1}}$, а $\sqrt[\leftroot{2}\uproot{4}m_{i+1}]{a_{i+1}}$ может переходить только в один из элементов вида $\sqrt[\leftroot{2}\uproot{4}m_{i+1}]{a_{i+1}}\xi_{m_{i+1}}^m$, где $m \in \{0, \dotsc, m_{i+1} - 1\}$, причем композиция автоморфизмов соответствует сложению степеней элемента $\xi_{m_{i+1}}$. Значит, $\Aut_{K_i}L_i \subset \Z_n$, поэтому группа $\Aut_{K_i}L_i$ "--- тоже абелева.
		
		Заметим теперь, что $\Aut_{K_i}L \normal \Aut_{L_{i + 1}}L$ и $\Aut_{L_{i}}L \normal \Aut_{K_i}L$, поскольку расширения $K_i \supset L_{i+1}$ и $L_i \supset K_i$ нормальны. Применим теорему о гомоморфизме к сопоставлениям $\phi \mapsto \phi|_{K_i}$ и $\psi \mapsto \psi|_{L_i}$, где $\phi \in \Aut_{L_{i + 1}}L$, $\psi \in \Aut_{K_i}L$, и получим, что $\Aut_{L_{i+1}}K_i \hm\cong \Aut_{L_{i + 1}}L / \Aut_{K_i}L$ и $\Aut_{K_i}L_i \cong \Aut_{K_i}L / \Aut_{L_i}L$. Следовательно, $\Gal{f} = \Aut_{L_n}{L} \varnormal \Aut_{K_{n-1}}L \varnormal \dotsb \varnormal \Aut_{K_0}L \varnormal \Aut_{L_0}L = \{\id\}$ "--- это цепочка нормальных подгрупп с абелевыми факторами.
		
\EndProof